\documentclass[11pt,a4paper]{article}
\usepackage[margin=1.7cm]{geometry}
\usepackage[T1]{fontenc}
\usepackage{lmodern}
\usepackage{microtype}
\usepackage{amsmath,amssymb,mathtools}
\usepackage{booktabs,array,tabularx,multicol}
\usepackage{enumitem}
\usepackage{xcolor}
\usepackage[most]{tcolorbox}
\usepackage{fancyhdr}
\usepackage{titlesec}
\usepackage{hyperref}
\usepackage{tikz}
\usepackage{setspace}
\usepackage{lastpage}

\definecolor{navy}{HTML}{17324D}
\definecolor{blue}{HTML}{2D6A8A}
\definecolor{teal}{HTML}{2A7F78}
\definecolor{gold}{HTML}{B77B16}
\definecolor{red}{HTML}{A23B3B}
\definecolor{pale}{HTML}{F5F8FB}
\definecolor{paleblue}{HTML}{EEF5F8}
\definecolor{paleteal}{HTML}{EEF8F6}
\definecolor{palegold}{HTML}{FFF8E8}
\definecolor{palered}{HTML}{FFF1F1}
\definecolor{darkgray}{HTML}{3E4650}

\hypersetup{colorlinks=true,linkcolor=blue,urlcolor=blue}
\setlength{\parindent}{0pt}
\setlength{\parskip}{5pt}
\setlist[itemize]{leftmargin=1.4em,itemsep=2pt,topsep=3pt}
\setlist[enumerate]{leftmargin=1.6em,itemsep=2pt,topsep=3pt}
\pagestyle{fancy}
\fancyhf{}
\lhead{\textcolor{navy}{\textbf{MATH6310 Metric Spaces \& Topology}}}
\rhead{\textcolor{darkgray}{Revision Guide}}
\cfoot{\textcolor{darkgray}{\thepage\,/\,\pageref{LastPage}}}
\renewcommand{\headrulewidth}{0.3pt}

\titleformat{\section}{\Large\bfseries\color{navy}}{\thesection}{0.6em}{}
\titleformat{\subsection}{\large\bfseries\color{blue}}{\thesubsection}{0.6em}{}
\titleformat{\subsubsection}{\normalsize\bfseries\color{teal}}{}{0em}{}

\newtcolorbox{keybox}[1][]{enhanced,breakable,colback=paleblue,colframe=blue,boxrule=0.7pt,arc=2mm,left=2.5mm,right=2.5mm,top=2mm,bottom=2mm,title=\textbf{Key idea},fonttitle=\bfseries,#1}
\newtcolorbox{defbox}[1][]{enhanced,breakable,colback=paleteal,colframe=teal,boxrule=0.7pt,arc=2mm,left=2.5mm,right=2.5mm,top=2mm,bottom=2mm,title=\textbf{Definition},fonttitle=\bfseries,#1}
\newtcolorbox{exambox}[1][]{enhanced,breakable,colback=palegold,colframe=gold,boxrule=0.7pt,arc=2mm,left=2.5mm,right=2.5mm,top=2mm,bottom=2mm,title=\textbf{Exam focus},fonttitle=\bfseries,#1}
\newtcolorbox{trapbox}[1][]{enhanced,breakable,colback=palered,colframe=red,boxrule=0.7pt,arc=2mm,left=2.5mm,right=2.5mm,top=2mm,bottom=2mm,title=\textbf{Common trap},fonttitle=\bfseries,#1}
\newtcolorbox{theorembox}[1][]{enhanced,breakable,colback=pale,colframe=navy,boxrule=0.8pt,arc=2mm,left=2.5mm,right=2.5mm,top=2mm,bottom=2mm,title=\textbf{Theorem / Proposition},fonttitle=\bfseries,#1}
\newcommand{\R}{\mathbb{R}}
\newcommand{\N}{\mathbb{N}}
\newcommand{\Z}{\mathbb{Z}}
\newcommand{\Q}{\mathbb{Q}}
\newcommand{\PP}{\mathcal{P}}
\newcommand{\Ran}{\operatorname{Ran}}
\newcommand{\id}{\operatorname{id}}

\begin{document}

\begin{titlepage}
\centering
\vspace*{2.2cm}
{\Huge\bfseries\color{navy} MATH6310\\[0.25cm] Metric Spaces \& Topology\par}
\vspace{0.6cm}
{\LARGE\color{blue} Sections 0--1 Simplified Revision Guide\par}
\vspace{1.0cm}
\begin{tcolorbox}[width=0.88\textwidth,colback=pale,colframe=navy,arc=3mm,boxrule=0.9pt]
\centering
\large
A concise, exam-friendly guide to \textbf{metric spaces, sets, functions, images, equivalence relations, countability and Cantor's theorem}.\\[0.4cm]
Built from the supplied lecture notes and exercises, with formal statements retained but explanations made easier to follow.
\end{tcolorbox}
\vfill
{\large \textbf{How to use this guide}}\\[0.25cm]
\begin{minipage}{0.82\textwidth}
\begin{itemize}
\item Read the \textcolor{blue}{Key idea} boxes first to understand the intuition.
\item Memorise the \textcolor{teal}{Definitions} precisely.
\item Use the \textcolor{gold}{Exam focus} boxes to practise standard proof structures.
\item Watch for the \textcolor{red}{Common trap} boxes: these are exactly where notation and logic often go wrong.
\end{itemize}
\end{minipage}
\vfill
{\small Based on MATH6310 Part I, Sections 0 and 1 lecture notes (07.ix.2026).}
\end{titlepage}

\tableofcontents
\newpage

\section{Section 0: What is a Metric Space?}

\begin{keybox}
A metric space is simply a set together with a rule for measuring distance. The power of the idea is that once a distance satisfies three basic axioms, we can study convergence, continuity and topology in many settings using the same language.
\end{keybox}

\subsection{The definition you must know}
\begin{defbox}
A \textbf{metric} on a set $S$ is a function
\[
d:S\times S\to [0,\infty)
\]
such that for all $x,y,z\in S$:
\begin{align*}
\text{(MS1)}\quad & d(x,y)=0 \iff x=y,\\
\text{(MS2)}\quad & d(x,y)=d(y,x),\\
\text{(MS3)}\quad & d(x,z)\le d(x,y)+d(y,z).
\end{align*}
The pair $(S,d)$ is called a \textbf{metric space}.
\end{defbox}

\begin{exambox}
When asked to prove that a proposed formula defines a metric, structure your answer as:
\begin{enumerate}
\item non-negativity / finiteness if needed;
\item $d(x,y)=0\iff x=y$;
\item symmetry;
\item triangle inequality.
\end{enumerate}
The triangle inequality is usually the only genuinely difficult part.
\end{exambox}

\subsection{The standard metric on $\R$}
\[
d(x,y)=|y-x|.
\]
This works because absolute value satisfies:
\[
|a|=0\iff a=0,\qquad |-a|=|a|,\qquad |a+b|\le |a|+|b|.
\]

\subsection{Euclidean metric on $\R^N$}
\[
d_2(x,y)=\sqrt{\sum_{i=1}^N (y_i-x_i)^2}=\|y-x\|_2.
\]
The triangle inequality comes from the Euclidean norm triangle inequality, itself linked to Cauchy--Schwarz.

\subsection{Discrete metric}
For any set $S$,
\[
d_{\mathrm{disc}}(s,t)=
\begin{cases}
0,&s=t,\\
1,&s\ne t.
\end{cases}
\]

\begin{keybox}
The same underlying set can carry very different metrics. For example, $\R$ with the usual metric and $\R$ with the discrete metric are different metric spaces.
\end{keybox}

\begin{exambox}
To prove the triangle inequality for the discrete metric, the only nontrivial case is $s\neq t$. Then $d(s,t)=1$, while $d(s,u)+d(u,t)\ge 1$ because at least one of $s\neq u$ or $u\neq t$ must hold.
\end{exambox}

\subsection{Counting metric on finite subsets}
For a finite set $S$ and $A,B\subseteq S$,
\[
d(A,B)=\#(A\triangle B),
\]
where
\[
A\triangle B=(A\setminus B)\cup(B\setminus A)
\]
is the symmetric difference.

\begin{keybox}
Interpretation: $d(A,B)$ counts how many elements you must ``toggle'' to turn $A$ into $B$.
\end{keybox}

\subsection{Metrics on spaces of functions}
Let $I=[0,1]$ and $f,g\in C(I;\R)$.

\subsubsection{Integration metric}
\[
d_1(f,g)=\int_0^1 |g(s)-f(s)|\,ds.
\]
This measures the \emph{total accumulated difference} between two functions.

\subsubsection{Supremum metric}
\[
d_{\sup}(f,g)=\sup_{s\in I}|g(s)-f(s)|.
\]
This measures the \emph{largest pointwise difference} between two functions.

\begin{trapbox}
For the integration metric, proving
\[
\int_0^1 |h|=0 \implies h=0
\]
uses continuity. For arbitrary functions this implication can fail in pointwise form.
\end{trapbox}

\subsection{Hausdorff metric: the idea}
For suitable nonempty, bounded, closed sets $K,J\subseteq \R^2$, the Hausdorff distance is the smallest amount $\delta$ by which each set must be ``fattened'' so that each contains the other.

\begin{keybox}
The Hausdorff metric measures distance between \textbf{sets themselves}, not between individual points.
\end{keybox}

\newpage
\section{Section 1: Sets and Basic Set Operations}

\subsection{Core notation}
\begin{center}
\begin{tabularx}{0.95\textwidth}{>{\bfseries}l X}
\toprule
Notation & Meaning\\
\midrule
$A\subseteq S$ & every element of $A$ lies in $S$\\
$A\subsetneq S$ & $A\subseteq S$ and $A\ne S$\\
$A\setminus B$ & elements in $A$ but not in $B$\\
$\PP(S)$ & the power set: all subsets of $S$\\
$A\triangle B$ & symmetric difference\\
\bottomrule
\end{tabularx}
\end{center}

For a family $\mathcal C$ of subsets of $S$,
\[
\bigcup \mathcal C=\{s\in S: \exists C\in\mathcal C,\ s\in C\},
\]
\[
\bigcap \mathcal C=\{s\in S: \forall C\in\mathcal C,\ s\in C\}.
\]

\subsection{De Morgan and distributive laws}
\[
A\cap\left(\bigcup \mathcal C\right)=\bigcup_{C\in\mathcal C}(A\cap C),
\]
\[
A\setminus\left(\bigcup \mathcal C\right)=\bigcap_{C\in\mathcal C}(A\setminus C),
\]
\[
A\setminus\left(\bigcap \mathcal C\right)=\bigcup_{C\in\mathcal C}(A\setminus C).
\]

\begin{exambox}
To prove equality of two sets, the standard method is:
\[
X=Y \quad\Longleftrightarrow\quad \forall s,\; s\in X \iff s\in Y.
\]
In practice, prove both inclusions $X\subseteq Y$ and $Y\subseteq X$.
\end{exambox}

\subsection{Empty families}
If $\mathcal C=\varnothing$, then
\[
\bigcup \mathcal C=\varnothing,
\qquad
\bigcap \mathcal C=S.
\]
The second identity depends on knowing the ambient universe $S$.

\subsection{Cartesian products}
\[
S\times T=\{(s,t):s\in S,\ t\in T\}.
\]
More generally,
\[
\prod_{\lambda\in\Lambda}S_\lambda
\]
is the set of indexed families $(s_\lambda)_{\lambda\in\Lambda}$ with $s_\lambda\in S_\lambda$.

\begin{trapbox}
If one factor is empty, the whole product is empty. If the index set $\Lambda$ itself is empty, the empty product contains one element by convention.
\end{trapbox}

\subsection{Axiom of Choice}
\begin{theorembox}
If $(S_\lambda)_{\lambda\in\Lambda}$ is a family of nonempty sets, then
\[
\prod_{\lambda\in\Lambda}S_\lambda\neq\varnothing.
\]
Equivalently, for a family of nonempty subsets $C_\lambda\subseteq S$, there exists a choice function $f:\Lambda\to S$ such that $f(\lambda)\in C_\lambda$ for every $\lambda$.
\end{theorembox}

\begin{keybox}
Think of the Axiom of Choice as saying: from every nonempty box in an arbitrary family of boxes, you may choose one object simultaneously.
\end{keybox}

\section{Functions}

\subsection{What a function actually contains}
A function $f:S\to T$ has three ingredients:
\begin{enumerate}
\item a domain $S$;
\item a codomain $T$;
\item a rule assigning exactly one value $f(s)\in T$ to every $s\in S$.
\end{enumerate}

\begin{trapbox}
A function is not determined only by its formula. Changing the domain or codomain can change whether it is injective, surjective or bijective.
\end{trapbox}

\subsection{Important examples}
\begin{itemize}
\item Identity: $\id_S(s)=s$.
\item Inclusion: $\iota_{A\to S}:A\to S$, $a\mapsto a$.
\item Indicator:
\[
\mathbf 1_A(s)=\begin{cases}1,&s\in A,\\0,&s\notin A.\end{cases}
\]
\item Constant function: $f(s)=t_0$ for all $s$.
\item Coordinate projection from a product: $(s_j)_{j\in J}\mapsto s_k$.
\end{itemize}

\subsection{Range and composition}
\[
\Ran f=\{f(s):s\in S\}\subseteq T.
\]
If $f:S\to T$ and $g:T\to U$, then
\[
(g\circ f)(s)=g(f(s)).
\]

\begin{trapbox}
Composition is generally \textbf{not commutative}: $f\circ g$ and $g\circ f$ may be different, or one of them may not even be defined.
\end{trapbox}

\subsection{Restriction, extension and co-restriction}
If $A\subseteq S$,
\[
f|_A:A\to T,\qquad a\mapsto f(a)
\]
is the restriction of $f$.

If $B\subseteq T$ and $\Ran f\subseteq B$, then viewing $f$ as a map $S\to B$ is a co-restriction.

\section{Injective, Surjective and Bijective Maps}

\begin{defbox}
For $f:S\to T$:
\begin{itemize}
\item \textbf{Injective}: $f(s_1)=f(s_2)\Rightarrow s_1=s_2$.
\item \textbf{Surjective}: $\Ran f=T$.
\item \textbf{Bijective}: both injective and surjective.
\end{itemize}
\end{defbox}

\begin{keybox}
\textbf{Injection:} no two inputs collide.\\
\textbf{Surjection:} no target is missed.\\
\textbf{Bijection:} perfect one-to-one matching.
\end{keybox}

\subsection{Composition facts}
For $f:S\to T$ and $g:T\to U$:
\begin{itemize}
\item if $f$ and $g$ are injective, then $g\circ f$ is injective;
\item if $g\circ f$ is injective, then $f$ is injective;
\item if $f$ and $g$ are surjective, then $g\circ f$ is surjective;
\item if $g\circ f$ is surjective, then $g$ is surjective;
\item if $f$ and $g$ are bijective, then $g\circ f$ is bijective.
\end{itemize}

\subsection{Left and right inverses}
\begin{theorembox}
Ignoring the special empty-set edge case noted in the lectures:
\begin{align*}
f\text{ injective} &\iff \exists g:T\to S\text{ with }g\circ f=\id_S,\\
f\text{ surjective} &\iff \exists h:T\to S\text{ with }f\circ h=\id_T,\\
f\text{ bijective} &\iff f\text{ has a two-sided inverse}.
\end{align*}
\end{theorembox}

\begin{exambox}
A standard way to prove injectivity is to start with
\[
f(x)=f(y)
\]
and algebraically deduce $x=y$. To prove surjectivity, start with an arbitrary target $t\in T$ and solve $f(s)=t$ for an input $s\in S$.
\end{exambox}

\section{Direct and Inverse Images}

\subsection{Direct image}
For $A\subseteq S$,
\[
f(A)=\{f(s):s\in A\}\subseteq T.
\]

\subsection{Inverse image}
For $C\subseteq T$,
\[
f^{-1}(C)=\{s\in S:f(s)\in C\}\subseteq S.
\]

\begin{trapbox}
The notation $f^{-1}(C)$ for an inverse image does \textbf{not} mean that $f$ is invertible. Inverse images are defined for every function.
\end{trapbox}

\subsection{What direct images preserve}
\[
f\left(\bigcup \mathcal A\right)=\bigcup_{A\in\mathcal A}f(A).
\]
But in general only
\[
f\left(\bigcap\mathcal A\right)\subseteq\bigcap_{A\in\mathcal A}f(A).
\]
Equality can fail when $f$ is not injective.

\subsection{What inverse images preserve}
Inverse images behave much more cleanly:
\[
f^{-1}\left(\bigcup\mathcal C\right)=\bigcup_{C\in\mathcal C}f^{-1}(C),
\]
\[
f^{-1}\left(\bigcap\mathcal C\right)=\bigcap_{C\in\mathcal C}f^{-1}(C),
\]
\[
f^{-1}(D\setminus C)=f^{-1}(D)\setminus f^{-1}(C).
\]

\begin{keybox}
Memory rule: inverse images preserve the usual set operations exactly. Direct images preserve unions, but intersections can become larger after mapping.
\end{keybox}

\section{Schr\"oder--Bernstein}
\begin{theorembox}
If there is an injection $S\to T$ and an injection $T\to S$, then there is a bijection $S\to T$.
\end{theorembox}

\begin{exambox}
The lecture notes explicitly state that you \textbf{must know the theorem}, but the proof is \textbf{not examinable}.
\end{exambox}

\section{Equivalence Relations and Partitions}

\begin{defbox}
A relation $\sim$ on a nonempty set $S$ is an \textbf{equivalence relation} if it is:
\begin{itemize}
\item reflexive: $s\sim s$;
\item symmetric: $s\sim t\Rightarrow t\sim s$;
\item transitive: $r\sim s$ and $s\sim t\Rightarrow r\sim t$.
\end{itemize}
\end{defbox}

For $a\in S$, the equivalence class is
\[
[a]=\{s\in S:s\sim a\}.
\]
The quotient set is
\[
S/\!\sim=\{[a]:a\in S\}.
\]

\subsection{Example: congruence modulo $p$}
For prime $p$,
\[
n\equiv_p m \iff p\mid(m-n).
\]
The equivalence classes are the residue classes
\[
[0]_p,[1]_p,\dots,[p-1]_p,
\]
and these split $\Z$ into disjoint pieces.

\subsection{Partitions}
\begin{defbox}
A collection $\mathcal C$ of subsets of $S$ is a partition if:
\begin{enumerate}
\item $\bigcup \mathcal C=S$;
\item every $C\in\mathcal C$ is nonempty;
\item distinct members of $\mathcal C$ are disjoint.
\end{enumerate}
\end{defbox}

\begin{theorembox}
Equivalence relations and partitions are two ways of describing the same structure:
\begin{itemize}
\item equivalence classes of an equivalence relation form a partition;
\item a partition defines an equivalence relation by saying two elements are related exactly when they lie in the same block.
\end{itemize}
\end{theorembox}

\begin{keybox}
An equivalence relation tells you \emph{when two elements count as the same}. A partition shows the resulting groups of ``same'' elements.
\end{keybox}

\subsection{Canonical decomposition of a function}
Given $f:S\to T$, define
\[
s_1\sim_f s_2 \iff f(s_1)=f(s_2).
\]
Then $S/\!\sim_f$ corresponds bijectively to $\Ran f$. Therefore every function decomposes canonically as
\[
S\twoheadrightarrow S/\!\sim_f\xrightarrow{\cong}\Ran f\hookrightarrow T.
\]

\begin{keybox}
Every function can be viewed as: first collapse together inputs with the same output, then make a perfect bijection with the range, then include the range inside the codomain.
\end{keybox}

\section{Countability}

\begin{defbox}
A set $S$ is \textbf{countable} if there exists an injection
\[
f:S\to\N.
\]
Thus finite sets are countable, and an infinite countable set can be listed in a sequence.
\end{defbox}

\subsection{Useful tests}
If $S\neq\varnothing$, then $S$ is countable iff there is a surjection $\N\to S$.

If $T$ is countable, then $S$ is countable whenever either:
\begin{itemize}
\item there is an injection $S\to T$, or
\item there is a surjection $T\to S$.
\end{itemize}

\subsection{Standard examples}
\begin{itemize}
\item $\N$ is countable.
\item $\Z$ is countable.
\item $\N^N$ is countable for each fixed finite $N$.
\item finite Cartesian products of countable sets are countable.
\item $\Q$ is countable because $\Z\times\N\to\Q$, $(n,m)\mapsto n/m$, is surjective.
\item a countable union of countable sets is countable.
\end{itemize}

\begin{exambox}
For countability questions, try one of three standard moves:
\begin{enumerate}
\item build an injection into a known countable set;
\item build a surjection from a known countable set;
\item write the target as a countable union of countable pieces.
\end{enumerate}
\end{exambox}

\subsection{Classic applications}
The lecture notes use the countable-union principle to show that:
\begin{itemize}
\item polynomials with rational coefficients form a countable set;
\item the set of all finite subsets of $\N$ is countable.
\end{itemize}

\section{Cantor's Theorem and Uncountability}

\begin{theorembox}
For every set $S$, there is no surjection
\[
S\to\PP(S).
\]
Equivalently, the power set $\PP(S)$ is strictly ``larger'' than $S$ in cardinality.
\end{theorembox}

\subsection{The diagonal idea}
Assume for contradiction that $\phi:S\to\PP(S)$ is surjective. Define
\[
D=\{s\in S:s\notin\phi(s)\}.
\]
Since $D\subseteq S$, surjectivity says $D=\phi(x)$ for some $x\in S$.

Now ask whether $x\in D$:
\begin{itemize}
\item If $x\in D$, then by definition $x\notin\phi(x)=D$.
\item If $x\notin D$, then by definition $x\in\phi(x)=D$.
\end{itemize}
Either way gives a contradiction.

\begin{keybox}
The diagonal construction creates an object deliberately designed to differ from the $s$th listed object at position $s$. Therefore no proposed list can contain everything.
\end{keybox}

\subsection{Binary sequences are uncountable}
The notes identify
\[
\PP(\N)\cong \{0,1\}^{\N}.
\]
By Cantor's theorem, these sets are uncountable.

\subsection{Uncountability of $\R$}
One exercise uses the injection
\[
\psi:\{0,1\}^{\N}\to\R,\qquad (a_n)\mapsto\sum_{n=1}^\infty\frac{a_n}{10^n}
\]
to deduce that $\R$ is uncountable.

\section{Optional / Non-examinable Material}

\subsection{Partial orders}
A partial order $\le$ satisfies reflexivity, antisymmetry and transitivity.

Important vocabulary:
\begin{itemize}
\item upper bound;
\item maximal element;
\item chain;
\item least upper bound.
\end{itemize}

\begin{exambox}
The lecture notes label partial orders and Zorn's Lemma as \textbf{self-study / non-examinable}. Learn the vocabulary if useful, but prioritise the examinable material first.
\end{exambox}

\newpage
\section{Proof Templates You Should Be Able to Reproduce}

\begin{enumerate}[label=\textbf{\arabic*.}]
\item \textbf{Set equality:} take an arbitrary element and prove membership both ways.
\item \textbf{Injectivity:} assume $f(x)=f(y)$ and derive $x=y$.
\item \textbf{Surjectivity:} let $t$ be arbitrary in the codomain and find $s$ with $f(s)=t$.
\item \textbf{Equivalence relation:} verify reflexive, symmetric, transitive separately.
\item \textbf{Metric:} verify zero-distance criterion, symmetry and triangle inequality.
\item \textbf{Countability:} construct an injection/surjection or decompose into countable pieces.
\item \textbf{Cantor diagonal argument:} assume a complete listing exists, construct an object that differs from every listed object on its diagonal position.
\end{enumerate}

\section{Fast Memory Sheet}

\begin{center}
\renewcommand{\arraystretch}{1.35}
\begin{tabularx}{\textwidth}{>{\bfseries\color{navy}}p{0.25\textwidth} X}
Metric & distance function obeying identity of indiscernibles, symmetry, triangle inequality\\
Injection & different inputs have different outputs\\
Surjection & every codomain element is hit\\
Bijection & injection + surjection\\
Direct image & push a subset of the domain forward\\
Inverse image & pull a subset of the codomain back\\
Equivalence relation & reflexive + symmetric + transitive\\
Partition & nonempty, pairwise disjoint blocks covering the whole set\\
Countable & injects into $\N$\\
Cantor & no surjection $S\to\PP(S)$\\
Schr\"oder--Bernstein & injections both ways imply a bijection\\
\end{tabularx}
\end{center}

\section{Mini Practice Check}
Try these without looking back.
\begin{enumerate}
\item State the three metric axioms.
\item Prove the discrete metric satisfies the triangle inequality.
\item Explain why the codomain matters when deciding surjectivity.
\item State the definitions of injective, surjective and bijective.
\item Why does $f^{-1}(C)$ make sense even if $f$ has no inverse function?
\item Give the three conditions for an equivalence relation.
\item Explain why equivalence classes form a partition.
\item Give one proof that $\Q$ is countable.
\item State Cantor's theorem and explain the diagonal set $D$.
\item Which parts of the supplied Section 1 notes are explicitly non-examinable?
\end{enumerate}

\begin{tcolorbox}[colback=navy!4,colframe=navy,arc=2mm]
\centering\large\bfseries\color{navy}
If you can state every definition precisely and reproduce the proof templates above, you have the core language needed for the next parts of Metric Spaces \& Topology.
\end{tcolorbox}

\end{document}
